%--------------------------------------------------------------------------------
%
%--------------------------------------------------------------------------------
\unnumberedChapter{RtLib Command Interface}

The Runtime Library provides two different ways to interact with an LCS node. During
normal operation, communication takes place through LCS messages exchanged over the
layout control bus. For firmware development, configuration and service purposes, the
runtime also provides a command interface.

The command interface exposes many of the internal runtime functions through a simple
text-based console. Any module providing a serial interface, for example through an USB
connector, can provide access to this console. A terminal program running on a computer
is sufficient to interact with the node.

The command interface is particularly useful during the development of new hardware and
firmware. It allows the programmer to inspect internal data structures, change
configuration values, simulate events and send messages without requiring additional
firmware code.

The command interface is not intended to replace the LCS message protocol. Instead, it
provides a local service interface to the runtime library.

%--------------------------------------------------------------------------------
\unnumberedSection{Serial Interface Overview}

The primary communication mechanism of the Layout Control System is the exchange of
LCS messages over the layout control bus. However, during development, configuration
and service operations, direct access to a node is extremely valuable.

The Runtime Library therefore provides a serial interface which can be accessed, for
example, through an USB connector on a hardware module. A standard terminal program is
sufficient to communicate with the node.

The serial interface provides two related but distinct functions:

\begin{smallGapItemize}
    \item a command interface for interacting with the runtime library,
    \item an ASCII representation of LCS messages for tools and communication gateways.
\end{smallGapItemize}

The command interface allows a user to inspect and modify node data, test event
processing and execute runtime functions. It is primarily intended for firmware
development, commissioning and troubleshooting.

The ASCII message representation provides a simple way to exchange LCS messages with
external tools. It can be used, for example, by a computer application or a gateway
which transports LCS messages over another communication medium such as Ethernet.

Neither interface replaces the normal LCS bus communication. Instead, both provide
additional access paths to the runtime and are especially useful during development and
service.


%--------------------------------------------------------------------------------
\unnumberedSection{ASCII Representation of LCS Messages}

In addition to processing command strings, the runtime library can convert LCS messages
between their binary representation and a simple ASCII text representation.

This format is useful for test programs, debugging tools and communication gateways.
A binary LCS message can be represented as a sequence of hexadecimal byte values
enclosed by the \texttt{<} and \texttt{>} delimiters.

\lstset{style=codesnippet-nobox}
\begin{lstlisting}
< data-byte-1 ... data-byte-n >
\end{lstlisting}

Each byte is represented by a hexadecimal value in the form \texttt{0xdd}, where
\texttt{d} represents a hexadecimal digit. A complete message therefore consists of
the opening delimiter, the individual message bytes and the closing delimiter.

For example:

\lstset{style=codesnippet-nobox}
\begin{lstlisting}
< 0x01 0x24 0x80 0x05 >
\end{lstlisting}

The Runtime Library provides helper functions to convert between the binary message
format used internally and the ASCII representation.

\lstset{style=api-interface}
\begin{lstlisting}
int lcsMsgToStr( char *msgStr,
                 uint8_t *dataBuf,
                 uint8_t dataBufLen );

int strToLcsMsg( uint8_t *dataBuf,
                 uint8_t dataBufLen,
                 char *msgStr );
\end{lstlisting}

The ASCII representation is deliberately simple. It is not intended as a replacement
for the LCS protocol, but as a convenient interchange format for software tools and
diagnostic applications.

\unnumberedSection{Command Processing}

When the runtime receives a command line, it first checks whether the command is
implemented by the core runtime library. The runtime provides commands for accessing
node data, managing configuration, inspecting internal state and testing message
processing.

Commands that are not recognized by the runtime are passed to an optional firmware
defined command callback. This allows a module designer to extend the command interface
with hardware or application specific commands.

The command interface therefore follows the same extension principle as the runtime
library itself:

\begin{itemize}
    \item the core runtime provides common functionality for all nodes,
    \item the firmware programmer can add application specific functionality.
\end{itemize}

The most important command when starting a session is the help command. It displays all
available commands together with their basic syntax.

\lstset{style=codesnippet-nobox}
\begin{lstlisting}
??? show example startup message
then type:

?

and display command list
\end{lstlisting}


%--------------------------------------------------------------------------------
\unnumberedSection{Configuration Mode Commands}

Many runtime operations modify persistent node data. To prevent accidental changes
during normal operation, these operations are restricted to configuration mode.

The following commands switch between configuration and operational mode.

\begin{table}[ht!]
    \begin{center}
        \renewcommand{\arraystretch}{1.2}
        \begin{tabularx}{0.8\textwidth}{|l|l|X|}
            \hline
            \textbf{Command} & \textbf{Arguments}  & \textbf{Operation} \\
            \hline
            C & - & switches node to configuration mode\\
            \hline
            O & - & switches node to operations mode\\
            \hline
        \end{tabularx}
    \end{center}
\end{table}
\FloatBarrier


When configuration mode is active, the command prompt indicates the current mode.

\lstset{style=codesnippet-nobox}
\begin{lstlisting}
C>          configuration mode
O>          operational mode
\end{lstlisting}

??? show prompt changes


%--------------------------------------------------------------------------------
\unnumberedSection{Accessing Node Data}

The command interface provides direct access to node and port data using the same
item-based concept introduced earlier in this book.

The commands correspond closely to the runtime library functions:

\begin{center}
\begin{tabular}{ll}
g & nodeGet() \\
p & nodeSet() \\
r & nodeReq()
\end{tabular}
\end{center}

The item number identifies the attribute or runtime function to access. This provides
a simple way to inspect and modify node data without writing dedicated firmware.

\begin{table}[ht!]
    \begin{center}
        \renewcommand{\arraystretch}{1.2}
        \begin{tabularx}{0.9\textwidth}{|l|l|X|}
            \hline
            \textbf{Command} & \textbf{Arguments}  & \textbf{Operation} \\
            \hline
            g & npId item [ val ] & retrieves an attribute \\
            \hline
            p & npId item val & sets an attribute \\
            \hline
            r & npId item [ val1 [ val2 ]] & request a function \\
            \hline
        \end{tabularx}
    \end{center}
\end{table}
\FloatBarrier



%--------------------------------------------------------------------------------
\unnumberedSection{Event Testing}

Events are a central concept in the LCS system. The command interface provides tools to
inspect and test event handling without requiring another node to generate the event.

Events can be locally generated, allowing the programmer to verify:

\begin{itemize}
    \item event configuration,
    \item event map entries,
    \item port reactions,
    \item callback execution.
\end{itemize}

\begin{table}[ht!]
    \begin{center}
        \renewcommand{\arraystretch}{1.2}
        \begin{tabularx}{0.9\textwidth}{|l|l|X|}
            \hline
            \textbf{Command} & \textbf{Arguments}  & \textbf{Operation} \\
            \hline
            e & mode nodeId eventId [ arg ] & simulate sending an event ( mode: 0 - on, 1 - off, 2 - evt ) \\
            \hline
        \end{tabularx}
    \end{center}
\end{table}
\FloatBarrier


Adding and removing event map entries is performed through runtime request functions
and therefore follows the same access rules as other persistent configuration data.


%--------------------------------------------------------------------------------
\unnumberedSection{Sending Raw LCS Messages}

For protocol development and debugging, the command interface provides a mechanism to
send raw messages directly onto the LCS bus.

This command intentionally bypasses the higher level runtime functions. It is useful
when developing new message types or verifying communication with other nodes.

\begin{table}[ht!]
    \begin{center}
        \renewcommand{\arraystretch}{1.2}
        \begin{tabularx}{0.9\textwidth}{|l|l|X|}
            \hline
            \textbf{Command} & \textbf{Arguments}  & \textbf{Operation} \\
            \hline
            B & byte1 [ ... byte8 ] & send a raw CAN bus data packet.\\
            \hline
        \end{tabularx}
    \end{center}
\end{table}
\FloatBarrier



%--------------------------------------------------------------------------------
\unnumberedSection{Runtime Diagnostics}

One of the most useful features of the command interface is the ability to inspect the
internal state of a node.

The \texttt{s} command provides formatted views of runtime data structures. This is
particularly valuable when debugging a node that does not behave as expected.

The command can display information from:

\begin{itemize}
    \item volatile memory (MEM),
    \item non-volatile storage (NVM),
    \item formatted runtime views.
\end{itemize}

\begin{table}[ht!]
    \begin{center}
        \renewcommand{\arraystretch}{1.2}
        \begin{tabularx}{0.9\textwidth}{|l|l|X|}
            \hline
            \textbf{Command} & \textbf{Arguments}  & \textbf{Operation} \\
            \hline
            s & sItem & display status view "sItem".\\
            \hline
        \end{tabularx}
    \end{center}
\end{table}
\FloatBarrier

The \texttt{s} command accepts a numeric parameter which indicates what status information to print. They are grouped for showing memory data, NVM data and a formatted version of some data.

\begin{table}[ht!]
    \begin{center}
        \renewcommand{\arraystretch}{1.2}
        \begin{tabularx}{0.9\textwidth}{|l|l|l|l|X|}
            \hline
            \textbf{Command} & \textbf{MEM}  & \textbf{NVM} & \textbf{FMT} & \textbf{Operation} \\
            \hline
            s & 0 & - & - & Board summary.\\
            \hline
            s & & 21 & - & Node Header.\\
            \hline
            s & 2 & 22 & 42 & Node Map.\\
            \hline
            s & 3 & 23 & - & Node Data.\\
            \hline
            s & - & 24 & - & Event Map.\\
            \hline
            s & 5 & 25 & 45 & Port Map.\\
            \hline
            s & 6 & - & 46 & Task Map.\\
            \hline
            s & - & - & 50 & I2C Devices.\\
            \hline
            s & - & - & 51 & Resource Descriptor Map.\\
            \hline
            s & - & - & 52 & Resource Map.\\
            \hline
        \end{tabularx}
    \end{center}
\end{table}
\FloatBarrier



The formatted views are intended for human inspection, while the MEM and NVM views are
useful when investigating storage or configuration problems.

% ??? examples


%--------------------------------------------------------------------------------
\unnumberedSection{Text Representation of LCS Messages}

In addition to commands, the serial interface can also be used to represent LCS messages
as ASCII text.

This representation is useful for communication gateways where messages need to be
transported through another medium, such as Ethernet. It also provides a convenient way
to create simple test tools.

The message is enclosed by the \texttt{<} and \texttt{>} delimiters and contains the
individual message bytes.

\lstset{style=codesnippet-nobox}
\begin{lstlisting}

< 0x01 0x02 0x03 >

\end{lstlisting}


The runtime library provides helper functions to convert between the binary message
representation and the ASCII format.

\lstset{style=api-interface}
\begin{lstlisting}

int lcsMsgToStr( char *msgStr,
                 uint8_t *dataBuf,
                 uint8_t dataBufLen );

int strToLcsMsg( uint8_t *dataBuf,
                 uint8_t dataBufLen,
                 char *msgStr );

\end{lstlisting}


%--------------------------------------------------------------------------------
\unnumberedSection{Summary}

The command interface provides a direct way to interact with an LCS node without
requiring additional firmware changes.

It exposes the same concepts used by the Runtime Library: item based data access,
runtime requests, event handling and message processing. This makes it a valuable tool
during firmware development, hardware bring-up and service operations.

A simple USB connection and a terminal program are sufficient to inspect a node, modify
configuration data and test communication.

The command interface is therefore not part of the normal operational communication
path. Instead, it provides a powerful service and diagnostic interface for developers
and users configuring an LCS system.





